\documentclass[11pt,a4paper]{article}
\usepackage[T1]{fontenc}
\usepackage[utf8]{inputenc}
\usepackage[czech]{babel}
\usepackage{lmodern,microtype}
\usepackage[left=30mm,right=20mm,top=24mm,bottom=25mm]{geometry}
\usepackage{amsmath,amssymb,booktabs,array,tabularx,longtable}
\usepackage{enumitem,xcolor,xurl}
\usepackage{fancyhdr}
\usepackage[unicode,colorlinks=true,linkcolor=black,citecolor=black,urlcolor=blue!45!black]{hyperref}
\setlength{\emergencystretch}{3em}
\setlength{\parskip}{0.35em}
\setlength{\parindent}{1.2em}
\setlist{itemsep=0.2em,topsep=0.4em,leftmargin=*}
\newcolumntype{Y}{>{\raggedright\arraybackslash}X}
\newcolumntype{L}[1]{>{\raggedright\arraybackslash}p{#1}}
\pagestyle{fancy}
\fancyhf{}
\fancyfoot[L]{\small Rešeršní studie \textbullet{} 21.~10.~2026}
\fancyfoot[R]{\small\thepage}
\renewcommand{\headrulewidth}{0pt}
\renewcommand{\footrulewidth}{0.3pt}
\raggedbottom
\newcommand{\studietitle}[3]{\noindent{\small\bfseries STUDIE #1 \hfill Verze 1.0 / 21. 10. 2026}\par\vspace{0.7em}
{\LARGE\bfseries #2\par}\vspace{0.5em}\noindent{\large #3\par}\vspace{0.8em}}
\newcommand{\zdroj}[3]{\bibitem{#1}#2. Dostupné z: \url{#3}. [cit.~2026-10-21].}
\begin{document}

\studietitle{01}{Přesnost RTK a podmínky jejího dosažení}{Princip metody a požadavky na stavbu měřicí sestavy}

\section{Zaměření a výchozí podmínky}
Studie vysvětluje, proč může metoda RTK dosahovat centimetrové přesnosti, které chyby tímto způsobem potlačit nelze a jaké podmínky musí splnit stavba měřicí sestavy pro vozidlo Formula Student. Výstupem je návrh montážních a ověřovacích kroků, nikoliv tvrzení o již dosažených výsledcích. Samostatné studie 02 a 03 řeší přenos RTCM a použití NTRIP.

Aplikačním východiskem jsou místní projektové podklady \emph{RTK GNSS a IMU pro Formula Student}, stav hlavního projektu \texttt{rtk-gnss-bakalarka} dostupný k datu rešerše: rover ArduSimple simpleRTK3B Pro se Septentrio mosaic-X5, pevná přenosná základna simpleRTK2B Budget se ZED-F9P, dvě antény ANN-MB2-00, Raspberry Pi 5 a zamýšlený spoj XIAO ESP32S3/Wio-SX1262. Jde o zvolenou architekturu, nikoli o potvrzení nákupu, osazení nebo funkčnosti. Přesná varianta F9P a verze firmwaru se musí při oživení zaznamenat.

Rešerše vychází z dokumentace výrobců a výukových podkladů ESA. Převzaté technické skutečnosti jsou citovány. Číselné příklady, návrhové meze a zkušební postupy jsou vlastní odvození pro tento projekt. Při přenosu závěrů do bakalářské práce je nutné zachovat toto rozlišení.

\section{V čem spočívá RTK}
\subsection{Kódové a fázové měření}
GNSS přijímač určuje vzdálenosti k družicím z přijatého rádiového signálu. Kódové měření využívá časový posun známé sekvence; fázové měření sleduje fázi nosné vlny. Fázi lze měřit podstatně jemněji než kód, avšak samotná fáze neříká, kolik celých vln mezi družicí a přijímačem leží. Tento neznámý celočíselný počet je fázová neurčitost neboli ambiguita. RTK využívá společná měření základny a roveru a odhad těchto neurčitostí.\cite{esa-rtk,esa-ambiguity}

Pro pochopení metody postačuje zjednodušená fázová rovnice v metrech
\begin{equation}
 \Phi_r^s=\rho_r^s+c(\delta t_r-\delta t^s)+T_r^s-I_r^s
          +\lambda N_r^s+\varepsilon_r^s .
\end{equation}
Index $r$ označuje přijímač a $s$ družici, $\rho$ geometrickou vzdálenost, $c$ rychlost světla, $\delta t$ chyby hodin, $T$ troposférické zpoždění, $I$ ionosférický člen, $\lambda$ vlnovou délku a $N$ ambiguitu. Zbytek $\varepsilon$ shrnuje šum a nevyjádřené vlivy. Záporné znaménko ionosférického členu platí pro fázi; kód má opačné. Rovnice je didakticky zjednodušená, bez úplného modelu přístrojových odchylek.

Rozdíl mezi dvěma přijímači a následně mezi dvěma družicemi vede k dvojité diferenci. Při vhodném modelu odstraní hodinové členy a dovolí řešit celočíselné ambiguity. Blízké přijímače navíc pozorují podobné atmosférické a orbitální chyby, jejichž rozdíl je menší než původní chyba. Tento princip nevylučuje zbytkové chyby ani špatné určení celého počtu vln. Přerušení sledování může způsobit skok o celý počet cyklů (cycle slip); příslušnou ambiguitu je pak nutné znovu vyřešit.\cite{esa-ambiguity}

\subsection{Základna, rover a režim FIX}
Pevná základna poskytuje svou referenční polohu a časově označené observace. Rover z nich spolu s vlastními měřeními určí vektor vůči základně. V běžném RTK tedy nejde jen o zasílání jedné univerzální opravy zeměpisné šířky a délky. Pro krátkou základnovou vzdálenost lze očekávat dobrou podobnost společných chyb; s rostoucí vzdáleností tato výhoda slábne. Zastínění a přerušení fázového sledování mohou vyvolat novou inicializaci.\cite{esa-rtk}

\par
\begin{table}[htbp]
\centering\small
\caption{Význam základních stavů řešení; žádný stav sám neurčuje chybu vůči nezávislé referenci.}
\begin{tabularx}{\linewidth}{L{27mm}YY}
\toprule
Stav & Co vyjadřuje & Interpretace při vyhodnocení\\
\midrule
Samostatné GNSS & Poloha bez řešení RTK vůči základně. & Nezahrnovat do souboru výsledků označených jako centimetrové RTK.\\
RTK FLOAT & Ambiguity jsou odhadovány jako reálné hodnoty. & Uvést zvlášť; přesnost závisí na podmínkách a době konvergence.\\
RTK FIX & Přijímač přijal celočíselné řešení ambiguit. & Nutná, avšak nepostačující informace pro důkaz správnosti polohy.\\
\bottomrule
\end{tabularx}
\end{table}
\par

FIX označuje výsledek rozhodnutí algoritmu. Nesprávná souřadnice základny, špatně zadaná výška antény či nesoulad souřadných rámců mohou existovat i při stabilním FIX. Proto tato studie navrhuje posuzovat stav řešení současně s odchylkou od reference, stářím korekcí, dostupností měření a kontinuitou záznamu.

\section{Co přesně znamená požadovaná přesnost}
\subsection{Relativní a absolutní poloha}
Vektorové vyjádření výsledku je
\begin{equation}
 \widehat{\mathbf x}_{R}=\widehat{\mathbf x}_{B}
                      +\widehat{\mathbf b}_{BR}.
\end{equation}
Z toho plyne důležitý konstrukční důsledek: chyba souřadnice základny se promítne do absolutní souřadnice roveru, i když je relativní vektor určen velmi přesně. Základna posunutá v zadání o jeden metr může vytvářet vzájemně přesné trajektorie posunuté v mapě přibližně o jeden metr. Přesná velikost přenosu chyby závisí na modelu a vzdálenosti; příklad vyjadřuje dominantní účinek chybného počátku.

Pro projekt je vhodné samostatně deklarovat:
\begin{itemize}
\item \textbf{přesnost relativní trajektorie} v místním systému tratě;
\item \textbf{opakovatelný souřadný počátek} pro porovnávání různých jízd a dnů;
\item \textbf{absolutní přesnost} vůči zaměřené referenci ve známém rámci a epoše.
\end{itemize}
Krátký autonomní \emph{survey-in} není automaticky geodetickým určením absolutní polohy. Pro opakovatelnost je nutné zachovat referenční bod a uložené souřadnice; při přestěhování antény je znovu určit. Manuál F9P rozlišuje režim survey-in a zadání pevné polohy základny.\cite{f9p}

\subsection{Přesnost, opakovatelnost a statistika}
Malý rozptyl bodů neprokazuje malou systematickou chybu. Výsledek se má popsat alespoň vodorovnou a svislou složkou, dobou měření, referencí a statistikou. Pro vodorovnou chybu vzorku $i$ lze použít
\begin{equation}
 e_{H,i}=\sqrt{(E_i-E_{\mathrm{ref},i})^2+
                   (N_i-N_{\mathrm{ref},i})^2}.
\end{equation}
Uvádět se má například medián, 95.~percentil, maximum a podíl platných vzorků. Výrobní údaj RMS, směrodatná odchylka a CEP nejsou zaměnitelné; bez znalosti rozdělení chyb nelze jeden přepsat na jiný.

Často používaný model přesnosti $a+b\,d$ kombinuje konstantní a vzdálenostní člen. Pro čistě ilustrační hodnoty $a=10$~mm a $b=1$~ppm vychází při $d=1$~km 11~mm, při $d=10$~km 20~mm. Tento výpočet \emph{není specifikací celé sestavy}. Neobsahuje chybu ustavení, rušení, pohyb držáku, referenčního bodu ani časové synchronizace.

\section{Požadavky na mechanickou a elektrickou stavbu}
\subsection{Anténa a její okolí}
Vícecestné šíření vzniká souběhem přímého a odraženého signálu. Protože je závislé na místním okolí antény, korekce z jiné lokality je obecně neodstraní. Významné jsou okolní odrazivé plochy, jejich geometrie a vlastnosti přijímací antény.\cite{multipath} Pro stavbu z toho vyplývá umístění s volným výhledem, odstup od velkých odrazivých ploch a zkouška přímo na vozidle, nikoliv pouze na laboratorním stole.

Fázový střed není totožný s libovolně zvoleným bodem plastového krytu. Rozlišují se mechanický referenční bod antény (ARP), offset fázového středu a jeho změny se směrem příjmu a frekvencí.\cite{antenna} Pro ANN-MB2 výrobce uvádí charakterizaci fázového středu na kruhové vodivé ploše o průměru 120~mm a upozorňuje na vliv změny této plochy.\cite{ann} Rozměr 120~mm proto představuje výchozí podmínku pro srovnání s dokumentací, nikoliv univerzální minimum pro každou montáž.

Navržené ověření montáže zahrnuje porovnání referenčního uspořádání s finálním držákem a krytem. Měnit se má vždy jedna vlastnost: podložka, kryt, výška nebo orientace. Zaznamená se přesná poloha ARP, průběhy kvality signálu, FIX a chyba polohy. Shodný typ antén na základně a roveru snižuje počet různých konstrukčních vlivů, ale sám nepotvrzuje shodnou kalibraci v odlišném prostředí.

U formule je třeba zkontrolovat pohled přes ochranný oblouk, pilota a karoserii. Vodivé či uhlíkové díly v okolí nelze posoudit jen podle vzhledu. Držák se nesmí při vibracích nebo aerodynamickém zatížení pohybovat vůči bodu vozidla, jehož trajektorie se vyhodnocuje. Kabel musí mít odlehčené konektory a opakovatelné vedení.

\subsection{Základna a ustavení}
Pro přenosnou základnu se navrhuje pevný stativ nebo jiný stabilní podklad mimo pohyb osob a vozidel. Zaznamená se označení bodu, výška a orientace antény, souřadnice, metoda jejich určení, souřadný rámec a datum. Přemístění antény při zachování starých souřadnic není platným pokračováním původního měření.

Pokud je cílem především srovnání jízd na jedné trati, lze nejprve pracovat s pevným místním počátkem a později jej navázat na přesné absolutní souřadnice. Tato volba musí být v datech jednoznačná. Výšku nad elipsoidem nelze bez převodu považovat za nadmořskou výšku. Při výměně korekční služby je třeba ověřit rámec a výškový systém, jinak mohou vzniknout zdánlivé skoky trajektorie.

\subsection{Napájení, rádiové rušení a rozhraní}
Konstrukční kontrola má pokrýt stabilitu napájení pod plnou zátěží, napájení aktivní antény podle konkrétní desky, úrovně UART a oddělení rušivých zdrojů. Hotová deska a samotný GNSS modul mají odlišné napájecí vstupy; závazné jsou návody konkrétních desek.\cite{board-f9p,board-x5}

Praktický A/B test se provede se zapnutým a vypnutým DC/DC měničem, Wi-Fi, LoRa, LTE modemem a ostatní elektronikou. Hodnotí se změna kvality příjmu a řešení, restarty i ztráty dat. Příčinou degradace může být vlastní zařízení, nikoliv družicová geometrie. Odolnost deklarovaná výrobcem GNSS modulu nenahrazuje tento test hotové sestavy.

\section{Konfigurace a návaznost na dynamiku vozidla}
Při oživení se ověří společné družicové signály, povolené RTCM zprávy a shodné ID základny. Vícepásmová anténa nerozšiřuje počet frekvencí, které skutečně měří základna. Pro korekci konkrétního signálu je rozhodující měření na obou stranách. Kompatibilitu je vhodné nejprve ověřit kabelem a teprve následně vložit rádiový spoj.

Rychlost výstupu polohy a frekvence přenosu korekcí jsou odlišné parametry. Samotné zvýšení výstupu roveru například na 50~Hz nevyžaduje automaticky 50~Hz RTCM, ale ani nevytváří 50 nezávislých potvrzení centimetrové přesnosti za sekundu. U projektu se má vyhodnotit latence polohy a její skutečné časové označení.

\subsection{Páka mezi anténou a sledovaným bodem}
Je-li $\mathbf l_{PA}^{B}$ vektor od sledovaného bodu $P$ k anténě $A$ v soustavě vozidla a $\mathbf R_B^E$ rotace do místní soustavy, platí
\begin{equation}
 \mathbf p_P^E=\mathbf p_A^E-\mathbf R_B^E\mathbf l_{PA}^{B}.
\end{equation}
Před použitím je třeba zjistit, zda výstup přijímače již obsahuje korekci anténního offsetu; stejná korekce se nesmí odečíst dvakrát. Pro pohybující se vozidlo nelze proměnnou páku nahrazovat konstantním posunem v ENU.

Vlastní orientační příklad: při délce páky 0,5~m a chybě orientace $1^\circ$ je příčná chyba přibližně $0{,}5\sin 1^\circ=8{,}7$~mm. Naklonění karoserie tedy může být významné i při dobrém GNSS řešení. Jedna anténa neposkytuje úplnou orientaci karoserie; směr pohybu se při smyku liší od jejího natočení.

\subsection{Časová synchronizace}
Přibližná chyba při časovém nesouladu $\Delta t$ mezi polohou a referencí je
\begin{equation}
 e_t \approx v|\Delta t|.
\end{equation}
Při modelové rychlosti 30~m/s znamená 1~ms 30~mm a 10~ms 300~mm. Jde o chybu přiřazení polohy k času, nikoliv o přímý vztah mezi stářím RTCM a chybou RTK. Stará korekce neznamená, že přijímač pouze vydává starou polohu vozidla.

PPS pomáhá navázat časové osy, avšak sám nezaručí, že každý vzorek IMU nebo každá zpráva na USB dostane přesný čas měření. Do ověření patří rozdíl epochy měření a příchodu zprávy, zpoždění ovladačů i chování Raspberry Pi pod zátěží. Čas přijmutí zprávy se nesmí bez ověření vydávat za čas měření.

\section{Navržené ověření dosažené přesnosti}
\subsection{Zkoušky po etapách}
\begin{enumerate}
\item \textbf{Kabelový základ.} Přímý přenos RTCM, otevřený výhled, stabilní antény. Ověřit konfiguraci a opakované získání FIX.
\item \textbf{Statická reference.} Nejméně tři nezávislé inicializace, více časových úseků s odlišnou geometrií. Porovnat s bodem zaměřeným metodou s vyčíslenou nejistotou.
\item \textbf{Finální mechanika.} Stejné zkoušky s držákem, krytem a napájením určeným do auta; zahrnout souběžný provoz všech vysílačů.
\item \textbf{Dynamická zkouška.} Porovnat anténní bod nebo přepočtený bod vozidla s nezávislou časově sladěnou referencí. Samotný návrat na stejnou stopu není referenční trajektorie.
\item \textbf{Poruchové stavy.} Výpadky korekcí, zastínění, restart rádia, znovupřipojení a změna základny. Sledovat chyby i dobu návratu do spolehlivého řešení.
\end{enumerate}
PPK z totožných observací může pomoci diagnostice, ale není plně nezávislou pravdou: sdílí antény, prostředí a část systematických chyb.

\subsection{Metriky a návrhová kritéria}
\par
\begin{table}[htbp]
\centering\small
\caption{Pracovní návrh přejímky; hodnoty nejsou dosažené výsledky ani výrobní garance.}
\begin{tabularx}{\linewidth}{L{40mm}YY}
\toprule
Ukazatel & Pracovní cíl & Podmínka interpretace\\
\midrule
95.~percentil vodorovné chyby & Do 30~mm ve vymezeném úseku s FIX. & Nezávislá reference a doložená časová nejistota.\\
95.~percentil svislé chyby & Do 50~mm při statické zkoušce. & Známý výškový systém a výška antény.\\
Dostupnost řešení & Alespoň 95\,\% času zvolené jízdy ve validovaném stavu. & Zvlášť uvést dobu inicializace a všechny výpadky.\\
Opakovatelnost ustavení & Rozdíl mezi sezeními vyhodnotit samostatně. & Stejný bod, rámec a montáž.\\
\bottomrule
\end{tabularx}
\end{table}
\par
Tyto cíle jsou návrhem pro další specifikaci experimentu. Jejich nesplnění se nesmí skrýt výběrem pouze nejlepších úseků. Vedle statistik pro FIX musí být zveřejněn výsledek za celou jízdu, nejdelší výpadek a počet chybných či neplatných řešení. Referenční nejistota má být dostatečně menší než ověřovaná mez; jinak závěr zkoušky zůstává neprůkazný.

Záznam má obsahovat konfiguraci a firmware, surová GNSS data obou přijímačů, přijaté korekce, časové značky, stav řešení, kvalitu signálu, počet použitých družic, odhad nejistoty a identifikaci základny. K tomu patří fotografie a rozměry montáže, podmínky zkoušky a známé zásahy do zařízení.

\section{Závěry pro projekt}
RTK může poskytnout centimetrové relativní měření, pouze pokud sestava zachová kvalitní fázové observace, správnou referenci a použitelná korekční data. Pro plánovanou formuli mají bezprostřední význam stabilní umístění obou antén, kontrola vodivé podložky ANN-MB2, rušení od vlastní elektroniky, správné souřadnice základny a časová návaznost GNSS/IMU.

První experiment má oddělit vlastnosti GNSS od komunikačního spoje: ověřit sestavu s přímým RTCM připojením a se známou referencí. Teprve potom je smysluplné hodnotit ztráty přesnosti způsobené LoRa nebo internetem. Za důkaz dosažené přesnosti se přijme kombinace reference, statistiky a dostupnosti, nikoli pouze dlouho svítící indikátor FIX.

\begin{thebibliography}{99}\small
\zdroj{esa-rtk}{GMV / ESA Navipedia. \emph{RTK Fundamentals}. 2011. Oddíly RTK Technique a Performance}{https://gssc.esa.int/navipedia/index.php/RTK_Fundamentals}
\zdroj{esa-ambiguity}{SANZ SUBIRANA, J.; JUAN ZORNOZA, J. M.; HERNÁNDEZ-PAJARES, M. \emph{Carrier Phase Ambiguity Fixing}. ESA Navipedia, 2011}{https://gssc.esa.int/navipedia/index.php/Carrier_Phase_Ambiguity_Fixing}
\zdroj{f9p}{U-BLOX. \emph{ZED-F9P Integration Manual}. UBX-18010802, R16, 30. 10. 2024. Oddíl 3.1.5, konfigurace RTK a základny}{https://content.u-blox.com/sites/default/files/ZED-F9P_IntegrationManual_UBX-18010802.pdf}
\zdroj{multipath}{SANZ SUBIRANA, J.; JUAN ZORNOZA, J. M.; HERNÁNDEZ-PAJARES, M. \emph{Multipath}. ESA Navipedia, 2011}{https://gssc.esa.int/navipedia/index.php/Multipath}
\zdroj{antenna}{SANZ SUBIRANA, J.; JUAN ZORNOZA, J. M.; HERNÁNDEZ-PAJARES, M. \emph{Antenna Phase Centre}. ESA Navipedia, 2011}{https://gssc.esa.int/navipedia/index.php/Antenna_Phase_Centre}
\zdroj{ann}{U-BLOX. \emph{ANN-MB2 Data Sheet}. UBXDOC-963802114-12775, R03, 2025. Oddíl 5.2, s. 9}{https://content.u-blox.com/sites/default/files/documents/ANN-MB2_DataSheet_UBXDOC-963802114-12775.pdf}
\zdroj{board-f9p}{ARDUSIMPLE. \emph{User Guide: simpleRTK2B Budget}}{https://www.ardusimple.com/user-guide-simplertk2b-budget/}
\zdroj{board-x5}{ARDUSIMPLE. \emph{User Guide: simpleRTK3B Pro}}{https://www.ardusimple.com/user-guide-simplertk3b-pro/}
\end{thebibliography}
\end{document}
